\begin{figure}[htbp]
\centering
% Setup: Skalierung der gesamten Zeichnung
\resizebox{.4836\textwidth}{!}{%
\begin{circuitikz}
% Allgemeiner Zeichnungsstil für alle Nodes
\tikzstyle{every node}=[font=\Huge]

% Adsorbens-Bereich (Tenax TA) im Inneren des Röhrchens
\draw [ color={rgb,255:red,92; green,92; blue,92} , fill={rgb,255:red,92; green,92; blue,92}, line width=1.5pt ] (22.5,-13.25) rectangle (42.5,-15.25);
% Quarzwolle-Stopfen (rechts/links) als Begrenzung des Adsorbens
\draw [ fill={rgb,255:red,0; green,0; blue,0} , line width=1.5pt ] (42.5,-13.25) rectangle (42.75,-15.25);
\draw [ fill={rgb,255:red,0; green,0; blue,0} , line width=1.5pt ] (22.5,-13.25) rectangle (22.25,-15.25);

% Haltefeder (schräg liegende Verbindung nahe der rechten Quarzwolle)
\draw [line width=1.5pt, short] (46.25,-13.5) -- (43,-13.5);
\draw [line width=1.5pt, short] (46.25,-15) -- (43,-15);
\draw [line width=1.5pt, short] (46.25,-13.5) -- (43,-15);

% Außenkontur des Glas-/Metallröhrchens (Rahmen)
\draw [line width=1.5pt, short] (52.5,-13) -- (52.5,-15.5);
\draw [line width=1.5pt, short] (52.5,-13) -- (22.5,-13);
\draw [line width=1.5pt, short] (52.5,-15.5) -- (22.5,-15.5);
% Kanten-Highlights der Außenkontur
\draw [ color={rgb,255:red,173; green,173; blue,173} , fill={rgb,255:red,173; green,173; blue,173}, line width=1.5pt ] (22.5,-13) rectangle (52.5,-13.25);
\draw [ color={rgb,255:red,173; green,173; blue,173} , fill={rgb,255:red,173; green,173; blue,173}, line width=1.5pt ] (22.5,-15.25) rectangle (52.5,-15.5);

% Messingkappe/Endkappe rechts (inkl. Zwischenstücke)
\draw [line width=1.5pt, short] (50,-13) -- (50,-11.75);
\draw [line width=1.5pt, short] (50,-15.5) -- (50,-16.5);
\draw [line width=1.5pt, short] (50,-16.25) -- (50,-16.75);
\draw [line width=1.5pt, short] (50,-11.75) -- (53.25,-11.75);
\draw [line width=1.5pt, short] (50,-16.75) -- (53,-16.75);
\draw [line width=1.5pt, short] (52.75,-16.75) -- (53.25,-16.75);
\draw [line width=1.5pt, short] (53.25,-11.75) -- (53.25,-12.5);
\draw [line width=1.5pt, short] (53.25,-16.75) -- (53.25,-16);
\draw [line width=1.5pt, short] (53.25,-12.5) -- (54,-12.5);
\draw [line width=1.5pt, short] (53.25,-16) -- (54,-16);
\draw [line width=1.5pt, short] (54,-12.5) -- (54,-11.75);
\draw [line width=1.5pt, short] (54,-16) -- (54,-16.75);
\draw [line width=1.5pt, short] (54,-11.75) -- (55,-11.75);
\draw [line width=1.5pt, short] (55,-11.75) -- (55,-16.75);
\draw [line width=1.5pt, short] (55,-16.75) -- (54,-16.75);
\draw [ color={rgb,255:red,214; green,214; blue,214} , fill={rgb,255:red,214; green,214; blue,214}, line width=1.5pt ] (50,-11.75) rectangle (53.25,-13);
\draw [ color={rgb,255:red,214; green,214; blue,214} , fill={rgb,255:red,214; green,214; blue,214}, line width=1.5pt ] (54,-11.75) rectangle (55,-16.75);
\draw [ color={rgb,255:red,214; green,214; blue,214} , fill={rgb,255:red,214; green,214; blue,214}, line width=1.5pt ] (53.25,-12.5) rectangle (54,-16);
\draw [ color={rgb,255:red,214; green,214; blue,214} , fill={rgb,255:red,214; green,214; blue,214}, line width=1.5pt ] (50.25,-15.5) rectangle (53.25,-16.75);
\draw [ color={rgb,255:red,214; green,214; blue,214} , fill={rgb,255:red,214; green,214; blue,214}, line width=1.5pt ] (50.25,-15.5) rectangle (50,-16.75);
\draw [ color={rgb,255:red,214; green,214; blue,214} , fill={rgb,255:red,214; green,214; blue,214}, line width=1.5pt ] (52.5,-13) rectangle (53.25,-15.5);
\draw [line width=1.5pt, short] (50,-11.75) -- (53.25,-11.75);
\draw [line width=1.5pt, short] (50,-16.75) -- (53.25,-16.75);
\draw [line width=1.5pt, short] (50,-15.5) -- (50,-16.75);
\draw [line width=1.5pt, short] (50,-11.75) -- (50,-13);
\draw [line width=1.5pt, short] (55,-11.75) -- (55,-16.75);
\draw [line width=1.5pt, short] (54,-16.75) -- (55,-16.75);
\draw [line width=1.5pt, short] (54,-11.75) -- (55,-11.75);
\draw [line width=1.5pt, short] (53.25,-12.5) -- (53.75,-12.5);
\draw [line width=1.5pt, short] (53.5,-12.5) -- (54,-12.5);
\draw [line width=1.5pt, short] (53.25,-16) -- (53.75,-16);
\draw [line width=1.5pt, short] (53.75,-16) -- (54,-16);
\draw [line width=1.5pt, short] (54,-16) -- (54,-16.75);
\draw [line width=1.5pt, short] (53.25,-16.75) -- (53.25,-16);
\draw [line width=1.5pt, short] (53.25,-11.75) -- (53.25,-12.5);
\draw [line width=1.5pt, short] (54,-12.5) -- (54,-11.75);
\draw [line width=1.5pt, short] (50,-13) -- (52.5,-13);
\draw [line width=1.5pt, short] (50,-15.5) -- (52.5,-15.5);
\draw [line width=1.5pt, short] (52.5,-13) -- (52.5,-15.5);
\draw [ color={rgb,255:red,173; green,173; blue,173} , fill={rgb,255:red,173; green,173; blue,173}, line width=1.5pt ] (22,-13.25) rectangle (22.5,-13.25);
\draw [ color={rgb,255:red,173; green,173; blue,173} , fill={rgb,255:red,173; green,173; blue,173}, line width=1.5pt ] (22,-15.25) rectangle (22.5,-15.25);
\draw [ color={rgb,255:red,133; green,133; blue,133}, line width=0.7pt, short] (22.5,-13) -- (22.25,-13.25);
\draw [ color={rgb,255:red,133; green,133; blue,133}, line width=0.7pt, short] (22.25,-13.25) -- (22,-13.25);
\draw [ color={rgb,255:red,133; green,133; blue,133}, line width=0.7pt, short] (22,-13.25) -- (21.75,-13);
\draw [ color={rgb,255:red,133; green,133; blue,133}, line width=0.7pt, short] (22.5,-15.5) -- (22.25,-15.25);
\draw [ color={rgb,255:red,133; green,133; blue,133}, line width=0.7pt, short] (22,-15.25) -- (21.75,-15.5);
\draw [ color={rgb,255:red,133; green,133; blue,133}, line width=0.7pt, short] (22,-15.25) -- (22.25,-15.25);
\draw [line width=0.7pt, short] (21.25,-13) -- (21.25,-13);
\draw [ color={rgb,255:red,173; green,173; blue,173} , fill={rgb,255:red,173; green,173; blue,173}, line width=1.5pt ] (22,-13.25) rectangle (21.75,-13.25);
\draw [ color={rgb,255:red,173; green,173; blue,173} , fill={rgb,255:red,173; green,173; blue,173}, line width=1.5pt ] (22,-15.25) rectangle (21.75,-15.25);
\draw [ color={rgb,255:red,173; green,173; blue,173} , fill={rgb,255:red,173; green,173; blue,173}, line width=1.5pt ] (21.75,-13.25) rectangle (18.75,-13);
\draw [ color={rgb,255:red,173; green,173; blue,173} , fill={rgb,255:red,173; green,173; blue,173}, line width=1.5pt ] (21.75,-15.25) rectangle (18.75,-15.5);
\draw [ color={rgb,255:red,173; green,173; blue,173} , fill={rgb,255:red,173; green,173; blue,173}, line width=1.5pt ] (18.75,-13.25) rectangle (17.5,-13);
\draw [ color={rgb,255:red,173; green,173; blue,173} , fill={rgb,255:red,173; green,173; blue,173}, line width=1.5pt ] (18.75,-15.5) rectangle (17.5,-15.25);
\draw [line width=1.5pt, short] (21.75,-13) -- (17.5,-13);
\draw [line width=1.5pt, short] (21.75,-15.5) -- (17.5,-15.5);
\draw [line width=1.5pt, short] (17.5,-13) -- (17.5,-15.5);
\draw [line width=1.5pt, short] (17.75,-13) -- (17.75,-15.5);

% DiffLok-Einlass (linkes Ende) mit Gehäuse und Öffnungen
\draw [ color={rgb,255:red,173; green,173; blue,173} , fill={rgb,255:red,173; green,173; blue,173}, line width=1.5pt ] (21.75,-13.25) rectangle (18.75,-13);
\draw [ color={rgb,255:red,173; green,173; blue,173} , fill={rgb,255:red,173; green,173; blue,173}, line width=1.5pt ] (21.75,-15.25) rectangle (18.75,-15.5);
\draw [line width=1.5pt, short] (21.75,-13) -- (17.5,-13);
\draw [line width=1.5pt, short] (21.75,-15.5) -- (17.5,-15.5);
\draw [line width=1.5pt, short] (17.5,-13) -- (17.5,-15.5);
\draw [line width=1.5pt, short] (17.75,-13) -- (17.75,-15.5);

% DiffLok-Filterstruktur (poröse Materialpunkte)
\draw [ color={rgb,255:red,194; green,194; blue,194} , fill={rgb,255:red,194; green,194; blue,194}, line width=0.2pt ] (19,-14.5) circle (0.15cm);
\draw [ color={rgb,255:red,194; green,194; blue,194} , fill={rgb,255:red,194; green,194; blue,194}, line width=0.2pt ] (20.15,-14) circle (0.15cm);
\draw [ color={rgb,255:red,194; green,194; blue,194}, fill={rgb,255:red,194; green,194; blue,194}, line width=0.2pt ] (19.5,-13.75) circle (0.15cm);
\draw [ color={rgb,255:red,194; green,194; blue,194}, fill={rgb,255:red,194; green,194; blue,194}, line width=0.2pt ] (20,-14.75) circle (0.15cm);
\draw [ color={rgb,255:red,194; green,194; blue,194}, fill={rgb,255:red,194; green,194; blue,194}, line width=0.2pt ] (21,-14.5) circle (0.15cm);
\draw [ color={rgb,255:red,194; green,194; blue,194}, fill={rgb,255:red,194; green,194; blue,194}, line width=0.2pt ] (18.25,-13.75) circle (0.15cm);
\draw [ color={rgb,255:red,194; green,194; blue,194}, fill={rgb,255:red,194; green,194; blue,194}, line width=1.5pt ] (25.5,-18) circle (0cm);
\draw [ color={rgb,255:red,194; green,194; blue,194}, fill={rgb,255:red,194; green,194; blue,194}, line width=0.2pt ] (21.75,-14) circle (0.15cm);

% Legendenpfeile zu Bauteilen (Beschriftungen)
\draw [line width=1.7pt, ->, >=Stealth] (54,-20.5) -- (54,-17)node[pos=0,below, fill=white]{Verschlusskappe};
\draw [line width=1.7pt, ->, >=Stealth] (42.5,-20.5) -- (42.5,-17)node[pos=0,below, fill=white]{Quarzwolle};
\draw [line width=1.7pt, ->, >=Stealth] (32.5,-20.5) -- (32.5,-17)node[pos=0,below, fill=white]{Adsorbens};
\draw [line width=1.7pt, ->, >=Stealth] (22.5,-20.5) -- (22.5,-17)node[pos=0,below, fill=white]{Quarzwolle};
\draw [line width=1.7pt, ->, >=Stealth] (17,-20.5) -- (17,-17)node[pos=0,below, fill=white]{Diffusionskappe};
\draw [line width=1.7pt, ->, >=Stealth] (48,-20.5) -- (45.5,-17)node[pos=0,below, fill=white]{Haltefeder};

% Strömungsrichtung des Trägergases (linke Pfeile)
\draw [line width=1.5pt, ->, >=Stealth] (13.25,-14.25) -- (16.75,-14.25);
\draw [line width=1.5pt, ->, >=Stealth] (13.25,-15) -- (16.75,-15);
\draw [line width=1.5pt, ->, >=Stealth] (13.25,-13.5) -- (16.75,-13.5);

% Anschluss- und Dichtungslinien nahe DiffLok
\draw [line width=1.5pt, short] (21.25,-13) -- (21.25,-12.25);
\draw [line width=1.5pt, short] (21.25,-15.5) -- (21.25,-16.25);
\draw [line width=1.5pt, short] (17.25,-13.25) -- (17.25,-12.25);
\draw [line width=1.5pt, short] (17.25,-15.25) -- (17.25,-16.25);
\draw [line width=1.5pt, short] (17.5,-16.5) -- (20.25,-16.5);
\draw [line width=1.5pt, short] (17.5,-12) -- (20,-12);
\draw [line width=1.5pt, short] (19.75,-12) -- (20.25,-12);
\draw [ color={rgb,255:red,173; green,173; blue,173}, line width=1pt, short] (17.25,-12.25) -- (17.5,-12);
\draw [ color={rgb,255:red,173; green,173; blue,173}, line width=1pt, short] (17.25,-16.25) -- (17.5,-16.5);
\draw [ color={rgb,255:red,173; green,173; blue,173}, line width=1pt, short] (21.25,-12.25) -- (21,-12);
\draw [ color={rgb,255:red,173; green,173; blue,173}, line width=1pt, short] (21.25,-16.25) -- (21,-16.5);
\draw [line width=1.5pt, short] (19.75,-12) -- (21,-12);
\draw [line width=1.5pt, short] (20,-16.5) -- (21,-16.5);
\draw [line width=1.5pt, ->, >=Stealth] (13.25,-13.5) -- (16.75,-13.5);
\draw [ color={rgb,255:red,214; green,214; blue,214} , fill={rgb,255:red,214; green,214; blue,214}, line width=0.2pt ] (8.75,-12.5) circle (0.15cm);
\draw [ color={rgb,255:red,214; green,214; blue,214} , fill={rgb,255:red,214; green,214; blue,214}, line width=0.2pt ] (8.5,-14.5) circle (0.15cm);
\draw [ color={rgb,255:red,214; green,214; blue,214} , fill={rgb,255:red,214; green,214; blue,214}, line width=0.2pt ] (10,-13.5) circle (0.15cm);
\draw [ color={rgb,255:red,214; green,214; blue,214}, fill={rgb,255:red,214; green,214; blue,214}, line width=0.2pt ] (9.5,-15.25) circle (0.15cm);
\draw [ color={rgb,255:red,214; green,214; blue,214}, fill={rgb,255:red,214; green,214; blue,214}, line width=0.2pt ] (10.75,-16) circle (0.15cm);
\draw [ color={rgb,255:red,214; green,214; blue,214}, fill={rgb,255:red,214; green,214; blue,214}, line width=0.2pt ] (8.75,-16.25) circle (0.15cm);
\draw [ color={rgb,255:red,214; green,214; blue,214}, fill={rgb,255:red,214; green,214; blue,214}, line width=0.2pt ] (10.75,-14.5) circle (0.15cm);
\draw [ color={rgb,255:red,214; green,214; blue,214}, fill={rgb,255:red,214; green,214; blue,214}, line width=0.2pt ] (11.75,-13.5) circle (0.15cm);
\draw [ color={rgb,255:red,214; green,214; blue,214}, fill={rgb,255:red,214; green,214; blue,214}, line width=0.2pt ] (12.5,-15) circle (0.15cm);
\draw [ color={rgb,255:red,214; green,214; blue,214}, fill={rgb,255:red,214; green,214; blue,214}, line width=0.2pt ] (11.25,-15.25) circle (0.15cm);
\draw [ color={rgb,255:red,214; green,214; blue,214}, fill={rgb,255:red,214; green,214; blue,214}, line width=0.2pt ] (10.75,-12.75) circle (0.15cm);
\draw [ color={rgb,255:red,214; green,214; blue,214}, fill={rgb,255:red,214; green,214; blue,214}, line width=0.2pt ] (12,-15.75) circle (0.15cm);
\draw [ color={rgb,255:red,214; green,214; blue,214}, fill={rgb,255:red,214; green,214; blue,214}, line width=0.2pt ] (12.5,-13) circle (0.15cm);
\draw [ color={rgb,255:red,214; green,214; blue,214}, fill={rgb,255:red,214; green,214; blue,214}, line width=0.2pt ] (11.75,-14.25) circle (0.15cm);
\draw [ color={rgb,255:red,214; green,214; blue,214}, fill={rgb,255:red,214; green,214; blue,214}, line width=0.2pt ] (9.25,-14.25) circle (0.15cm);
\draw [ color={rgb,255:red,214; green,214; blue,214}, fill={rgb,255:red,214; green,214; blue,214}, line width=0.2pt ] (11.25,-16.75) circle (0.15cm);
\draw [ color={rgb,255:red,214; green,214; blue,214}, fill={rgb,255:red,214; green,214; blue,214}, line width=0.2pt ] (9.75,-16) circle (0.15cm);
\draw [ color={rgb,255:red,214; green,214; blue,214}, fill={rgb,255:red,214; green,214; blue,214}, line width=0.2pt ] (9.75,-11.75) circle (0.15cm);
\draw [line width=1.5pt, ->, >=Stealth] (10,-20.5) -- (10,-17)node[pos=0,below, fill=white]{Probenmatrix};

% Beschriftung im Innenraum
\node [font=\Huge] at (15,-12.5) {Trägergas};

% Maßpfeile für Länge und Durchmesser
%\draw [line width=1.5pt, <->, >=Stealth] (17.5,-9.25) -- (52.5,-9.25)node[pos=0.5,above, fill=white]{8 cm};
%\draw [line width=1.5pt, <->, >=Stealth] (57.5,-13) -- (57.5,-15.5)node[pos=0.5,right, fill=white]{4 mm};

\end{circuitikz}
}
% Abbildungsunterschrift und Label
\caption{Schematische Darstellung eines Tenax Adsorbtionsröhrchens. Angepasst aus Lit. \cite{environments4040087}.}
\label{fig:tenax-adsorptionstube}
\end{figure}